\documentclass{article}

\usepackage[utf8]{inputenc}
\usepackage[T1]{fontenc}
\usepackage{amsmath,amssymb,amsthm}
\usepackage{mathtools}
\usepackage{hyperref}
\usepackage{booktabs}

\makeatletter
\newcommand{\citep}[2][]{%
  \begingroup
  \def\@tmpkey{#2}%
  \ifx\@tmpkey\@empty\else
    (\ifx\@empty#1\@empty\else #1, \fi
     \@nameuse{bibkey@\@tmpkey})%
  \fi
  \endgroup
}
\expandafter\def\csname bibkey@lasser2026exo\endcsname{Exogenous~Verification}
\expandafter\def\csname bibkey@lasser2026phase\endcsname{Phase~Redundancy}
\expandafter\def\csname bibkey@lasser2026micro\endcsname{Microfoundation}
\makeatother

\newtheorem{definition}{Definition}
\newtheorem{lemma}{Lemma}
\newtheorem{theorem}{Theorem}
\newtheorem{remark}{Remark}
\newtheorem{assumption}{Assumption}

\newcommand{\Aadv}{A_{\mathrm{adv}}}
\newcommand{\deltaadv}{\delta_{\mathrm{adv}}}
\newcommand{\deltastar}{\delta_{*}}
\newcommand{\rhogap}{\rho_{\mathrm{gap}}}
\newcommand{\Pproxy}{P}
\newcommand{\Ttruth}{T}
\newcommand{\epsgap}{\varepsilon_{\mathrm{gap}}}
\newcommand{\epsnonres}{\varepsilon_{\mathrm{gap}}^{\mathrm{nonres}}}
\newcommand{\epsfloor}{\varepsilon_{\mathrm{floor}}^{\mathrm{res}}}

\title{Lemma 6 ($h_{\mathrm{detect}}$ intensivity): Full Proof Draft v2}
\author{Paper 10 working draft for codex re-review}
\date{}

\begin{document}

\maketitle

\section{Goal and changes from v1}

V1 had two structural problems:
(a)~$h_{\mathrm{detect}}$ defined as integrated-increment $\int (d\epsgap/dt) dt$,
which bounds the gap \emph{change} during the detection window but
not the point-in-time gap $\epsgap(t)$ that Theorem~1's Layer~2
needs;
(b)~the per-unit-time formulation didn't specify the monitoring
clock rigorously.

V2 fixes both, plus refinements:

\textbf{Changes from v1:}
\begin{itemize}
\item \textbf{Q1 fix.}  (C11) stated in monitoring-clock form, with
  discrete (per-SPRT-step) and continuous variants explicitly
  noted.  The SPRT exposure clock $s$ is the canonical clock for
  per-step rate bounds.
\item \textbf{Q2 fix (the big one).}  $h_{\mathrm{detect}}$
  redefined as the supremum gap during the detection window, with
  the baseline term explicitly bounded:
  \[
    h_{\mathrm{detect}}(\theta) \;:=\; \sup_{t \in
    [t_0, t_0+T_{\mathrm{detect}}]} \epsgap(t)
    \;\leq\; \epsgap(t_0) + \rhogap(\deltastar) \cdot
    T_{\mathrm{detect}}.
  \]
  At the Layer~1$\to$Layer~2 boundary, $\epsgap(t_0) \leq
  h_{\mathrm{static}}(\theta)$ (the Layer~1 bound, intensive by
  Lemma~1), making $h_{\mathrm{detect}}$ intensive as a sum of two
  intensive quantities.
\item \textbf{Q3 fix.}  Renamed $\rho_g \to \rhogap$ to make
  explicit that (C11) bounds the gap-growth rate, not the
  alignment-property-error rate.  No double-counting with
  $K_{\mathrm{Lip}}$.
\item \textbf{Q4 fix.}  Use union-class KL floor $\deltastar =
  \min(\deltaadv, \delta_{\mathrm{adv}}^{\mathrm{env}})$ since
  Layer~2 covers $\Aadv \cup \Aadv^{\mathrm{env}}$.
\item \textbf{Q5 fix.}  (C11) explicitly framed for promotion to
  theorem-level condition.
\item \textbf{Source-citation correction.}
  \citep[Microfoundation]{lasser2026micro}'s ``$T$ failure modes''
  remark does not establish the per-step gap-growth rate; (C11) is
  treated as a new Paper~10 operational assumption.
\end{itemize}

\section{Setup: the monitoring clock and $h_{\mathrm{detect}}$}

Theorem~\ref{thm:deployment_safety}'s Layer~2 covers adversarial
strategies $q \in \Aadv \cup \Aadv^{\mathrm{env}}$ during the
detection window
$[t_0, t_0 + T_{\mathrm{detect}}]$, where $t_0$ is the violation
onset and $T_{\mathrm{detect}}$ is the SPRT detection time
(Lemma~\ref{lem:4}).

Layer~2 applies $h_{\mathrm{detect}}$ as a point-in-time bound:
\[
  |g(\Ttruth) - g(\Pproxy)|(t) \;\leq\; \mathrm{Lip}(g) \cdot
  h_{\mathrm{detect}}(\theta)
  \quad \text{for all } t \in [t_0, t_0 + T_{\mathrm{detect}}].
\]
Therefore $h_{\mathrm{detect}}$ must bound $\epsgap(t)$ at every
$t$ in the window:
\[
  h_{\mathrm{detect}}(\theta) \;:=\; \sup_{t \in
  [t_0, t_0 + T_{\mathrm{detect}}]} \epsgap(t).
\]

\paragraph{The monitoring clock.}
The SPRT machinery (Paper~5) operates on a per-event clock: each
ledger entry triggers an SPRT update, and the per-step KL
divergence drives detection.  Wall-clock time is related to SPRT
clock by the event-rate of the deployment.  For Lemma~6's purposes,
we use the SPRT exposure clock $s$ counted in events; converting
to wall-clock time requires multiplying by the average inter-event
interval.

For the deployment-safety theorem, the relevant detection-window
length is $T_{\mathrm{detect}}$ measured in SPRT steps (or,
equivalently, in events processed by the monitor).  The gap
growth rate $\rhogap$ is per-SPRT-step.

\section{Bounded gap-growth-rate condition (C11)}

\begin{assumption}[(C11): Bounded gap-growth rate per SPRT step]
\label{ass:C11}
For any adversarial strategy $q \in \Aadv \cup
\Aadv^{\mathrm{env}}$ producing channel deviations bounded by
$\deltastar = \min(\deltaadv, \delta_{\mathrm{adv}}^{\mathrm{env}})$
(the Lemmas~5b and~5e KL floors), the per-SPRT-step proxy-truth
gap growth is bounded:
\[
  (\Delta \epsgap)_n \;:=\; \epsgap(s_n) - \epsgap(s_{n-1}) \;\leq\;
  \rhogap(\deltastar),
\]
where $\rhogap(\deltastar)$ is a deployment-class constant
depending on $\mathcal{D}$'s operational parameters (channel
sensitivity-to-gap mapping, event-rate normalization) but
independent of $|P|$.

\textbf{Continuous-time variant.}  Equivalently, in continuous
SPRT exposure clock $s$:
\[
  \frac{d \epsgap}{ds} \;\leq\; \rhogap(\deltastar).
\]
The discrete and continuous forms agree under fixed step duration
or bounded event-rate normalization.
\end{assumption}

\textbf{Operational interpretation.}  Channel deviations bounded by
$\deltastar$ (else SPRT fires faster) limit how much $\epsgap$ can
grow per SPRT step.  The constant $\rhogap$ characterizes the
deployment's channel-deviation-to-gap-growth coupling.

\textbf{(C11) is new content.}  \citep[Microfoundation]{lasser2026micro}'s
``$T$ failure modes'' remark notes adequacy/failure modes of the
operational truth $T$ but does not establish a formal per-step
gap-growth-rate bound.  (C11) is a Paper~10 operational assumption,
analogous to (C7) bounded co-evolution and SA1 HHI surrogate
adequacy: deployment-class condition with empirical testability,
not a derived consequence of source-paper machinery.

\section{Statement of Lemma 6 (v2)}

\begin{lemma}[$h_{\mathrm{detect}}$ intensivity]
\label{lem:6}
Under Theorem~\ref{thm:deployment_safety}'s conditions (C5)
continuous SPRT monitoring and (C11) bounded gap-growth rate,
combined with Lemma~\ref{lem:4} (SPRT lead-time tail bound) and
Lemma~\ref{lem:5d} (lead-time tail composition):

The detection-window supremum proxy-truth gap satisfies:
\[
  h_{\mathrm{detect}}(\theta) \;\leq\; \epsgap(t_0) +
  \rhogap(\deltastar) \cdot T_{\mathrm{detect}},
\]
where $\epsgap(t_0)$ is the gap at violation onset.

At the Layer~1$\to$Layer~2 boundary (the moment when an adversarial
strategy first crosses the Layer~1 safe region), $\epsgap(t_0) \leq
h_{\mathrm{static}}(\theta)$ by Layer~1's intensivity argument.
Therefore:
\[
  h_{\mathrm{detect}}(\theta) \;\leq\; h_{\mathrm{static}}(\theta)
  + \rhogap(\deltastar) \cdot T_{\mathrm{detect}}.
\]

With probability at least $1 - \beta'$ (Lemma~\ref{lem:5d}'s
tail bound), $T_{\mathrm{detect}} \leq T_{\mathrm{cascade}}$, so:
\[
  h_{\mathrm{detect}}(\theta) \;\leq\; h_{\mathrm{static}}(\theta)
  + \rhogap(\deltastar) \cdot T_{\mathrm{cascade}}.
\]

$h_{\mathrm{detect}}(\theta)$ is intensive in $|P|$ over the
deployment class $\mathcal{D}$: $h_{\mathrm{static}}(\theta)$ is
intensive by Lemma~\ref{lem:1}, $\rhogap(\deltastar)$ is intensive
by (C11), and $T_{\mathrm{cascade}}$ is intensive by
\citep[Phase~Redundancy]{lasser2026phase}'s metastable-lifetime
bound.
\end{lemma}

\section{Proof}

\begin{proof}
We bound $h_{\mathrm{detect}}$ as the supremum of $\epsgap$ over
the detection window, then apply (C11) and the tail bound.

\emph{Step 1 (per-step decomposition).}  For any SPRT-step
sequence $s_0 = 0, s_1, \ldots, s_{N_{\mathrm{detect}}}$ during the
detection window:
\[
  \epsgap(s_n) \;=\; \epsgap(s_0) + \sum_{i=1}^{n}
  (\Delta \epsgap)_i.
\]

\emph{Step 2 (apply (C11)).}  By Assumption~\ref{ass:C11}, each
per-step increment satisfies $(\Delta \epsgap)_n \leq
\rhogap(\deltastar)$.  Therefore:
\[
  \epsgap(s_n) \;\leq\; \epsgap(s_0) + n \cdot
  \rhogap(\deltastar).
\]

\emph{Step 3 (supremum over window).}  The supremum over the
detection window of length $N_{\mathrm{detect}}$ SPRT steps:
\[
  h_{\mathrm{detect}}(\theta) \;=\; \sup_{0 \leq n \leq
  N_{\mathrm{detect}}} \epsgap(s_n)
  \;\leq\; \epsgap(s_0) + N_{\mathrm{detect}} \cdot
  \rhogap(\deltastar).
\]
Identifying $T_{\mathrm{detect}}$ with the number of SPRT steps to
detection (under fixed-step normalization):
\[
  h_{\mathrm{detect}}(\theta) \;\leq\; \epsgap(t_0) +
  \rhogap(\deltastar) \cdot T_{\mathrm{detect}}.
\]

\emph{Step 4 (boundary condition on $\epsgap(t_0)$).}  At the
Layer~1$\to$Layer~2 boundary, the adversarial strategy has just
crossed the Layer~1 safe region.  Immediately before the crossing,
Layer~1's bound holds:
$\epsgap(t_0^-) \leq h_{\mathrm{static}}(\theta)$ (by Layer~1
proof + Lemma~\ref{lem:1}).
By continuity (or single-step boundedness under (C11)), the gap at
$t_0$ is bounded by the same intensive constant within
$\rhogap(\deltastar) \cdot \delta s$ for one SPRT step;
absorbing this into the bound:
\[
  \epsgap(t_0) \;\leq\; h_{\mathrm{static}}(\theta).
\]

\emph{Step 5 (combine).}  Substituting Step~4 into Step~3:
\[
  h_{\mathrm{detect}}(\theta) \;\leq\; h_{\mathrm{static}}(\theta)
  + \rhogap(\deltastar) \cdot T_{\mathrm{detect}}.
\]

\emph{Step 6 (apply Lemma 5d tail bound).}  By Lemma~\ref{lem:5d},
$\Pr[T_{\mathrm{detect}} > T_{\mathrm{cascade}}] \leq \beta' \leq
\exp(-\kappa \tau_{\mathrm{meta}} \deltastar)$.  On the event
$T_{\mathrm{detect}} \leq T_{\mathrm{cascade}}$ (probability at
least $1 - \beta'$):
\[
  h_{\mathrm{detect}}(\theta) \;\leq\; h_{\mathrm{static}}(\theta)
  + \rhogap(\deltastar) \cdot T_{\mathrm{cascade}}.
\]

\emph{Step 7 (intensivity).}  Each constant is independent of
$|P|$ over $\mathcal{D}$:
\begin{itemize}
\item $h_{\mathrm{static}}(\theta)$: intensive by Lemma~\ref{lem:1}.
\item $\rhogap(\deltastar)$: deployment-class constant by (C11).
\item $T_{\mathrm{cascade}} \geq \tau_{\mathrm{meta}}$:
  \citep[Phase~Redundancy]{lasser2026phase}'s metastable-lifetime
  bound, deployment-class.
\end{itemize}
Therefore $h_{\mathrm{detect}}(\theta)$ is intensive in $|P|$ over
$\mathcal{D}$.  $\Box$
\end{proof}

\section{Discussion}

\paragraph{Why Lemma~5d's tail bound, not Lemma~4 directly.}
Lemma~5d composes Lemma~5b's KL floor with Lemma~4's tail bound,
giving the operational lead-time guarantee $\beta' \leq \exp(-\kappa
\tau_{\mathrm{meta}} \deltaadv)$.  V1 used $\deltaadv$ (Lemma~5b
only); v2 uses $\deltastar = \min(\deltaadv,
\delta_{\mathrm{adv}}^{\mathrm{env}})$ to cover Layer~2's full
detection class $\Aadv \cup \Aadv^{\mathrm{env}}$.  This requires
Lemma~5d to be re-stated with the union-class floor — a small
amendment that should be made at integration.

\paragraph{Constant taxonomy.}
\begin{itemize}
\item $h_{\mathrm{static}}(\theta)$: intensive composition
  (Lemma~\ref{lem:1}).
\item $\rhogap(\deltastar)$: \emph{newly assumed Paper 10
  operational constant} via (C11); calibrated by per-channel
  sensitivity-to-gap measurement.
\item $T_{\mathrm{cascade}}, \tau_{\mathrm{meta}}$: source-derived
  from \citep[Phase~Redundancy]{lasser2026phase} (with the (C7)
  bounded co-evolution caveat).
\item $\beta'$: composed from Lemmas~5b/5e, 4, 5d.
\end{itemize}

\paragraph{Layer 2 operational form.}
Combining Lemma~6 with $\mathrm{Lip}(g) \leq K_{\mathrm{Lip}}$ from
(C6):
\[
  |g(\Ttruth) - g(\Pproxy)|(t) \;\leq\; K_{\mathrm{Lip}} \cdot
  \bigl(h_{\mathrm{static}}(\theta) + \rhogap(\deltastar) \cdot
  T_{\mathrm{cascade}}\bigr)
\]
with probability at least $1 - \beta'$.  The right-hand side is the
sum of intensive constants times $K_{\mathrm{Lip}}$, all $|P|$-
independent.

\section{What this proof does and does not establish}

\textbf{Establishes:}
\begin{itemize}
\item Formal monitoring-clock formulation of (C11) with discrete
  (per-SPRT-step) and continuous variants.
\item Supremum-form $h_{\mathrm{detect}}$ definition matching
  Theorem 1's Layer 2 application.
\item Boundary-condition argument bounding $\epsgap(t_0)$ via
  Layer 1's $h_{\mathrm{static}}$.
\item Union-class KL floor $\deltastar$ for Layer 2's full
  detection class.
\item Intensive bound on $h_{\mathrm{detect}}$ as $h_{\mathrm{static}}
  + \rhogap T_{\mathrm{cascade}}$ with high probability.
\end{itemize}

\textbf{Does not establish (deferred):}
\begin{itemize}
\item Tightness of $\rhogap(\deltastar)$.  Operators may calibrate
  tighter bounds for restricted threat models.
\item Existence of deployment configurations satisfying (C11).
  Verified empirically (a new Audit~7).
\item Numerical value of $\rhogap(\deltastar)$.  Calibrated
  deployment-by-deployment.
\item Structural derivation of (C11) from finer source-paper
  machinery.  Currently a Paper 10 operational assumption.
\end{itemize}

\section{Specific verification questions for v2 codex review}

\begin{enumerate}
\item Is the supremum-form $h_{\mathrm{detect}}$ definition
  correct, and does the boundary condition $\epsgap(t_0) \leq
  h_{\mathrm{static}}(\theta)$ hold rigorously?  In particular, is
  the ``immediately before crossing'' continuity argument
  watertight, or does it require additional structure (e.g.,
  bounded jump at the boundary crossing)?

\item Is the SPRT-step monitoring-clock formulation of (C11)
  correctly stated?  The discrete/continuous variants agree under
  fixed-step normalization; is this normalization a separate
  assumption that should be named, or implicit in (C5) continuous
  SPRT monitoring?

\item Is the union-class KL floor $\deltastar = \min(\deltaadv,
  \delta_{\mathrm{adv}}^{\mathrm{env}})$ the right composition for
  Layer 2?  Lemma 5d as currently stated covers $\Aadv$ only;
  does the integration need to amend Lemma 5d to use $\deltastar$,
  or can Lemma 6 invoke Lemma 5b/5e separately and take the
  minimum?

\item Is the renaming $\rho_g \to \rhogap$ sufficient to clarify
  that (C11) bounds gap-growth rate, not alignment-property-error
  rate?  Is there any remaining double-counting risk with
  $K_{\mathrm{Lip}}$ in the Layer 2 operational form?

\item Should (C11) be (C11) of Theorem 1 (next in sequence) or
  could it be absorbed into (C5) continuous SPRT monitoring as a
  sub-clause?  The trade-off: separate condition makes the
  assumption visible; sub-clause keeps the condition list shorter.
\end{enumerate}

\end{document}
